\section{Algorithm Family and Effective Penalty}
\label{sec:algorithm-framework}

\subsection{Algorithm family and worst-case iteration complexity}
\label{subsec:algorithm-family}

For \(\rho>0\), the reduced augmented Lagrangian of~\eqref{prob:split} is
\begin{equation}
    \calD_\rho(\bx,\blmbd)
    :=\min_{\by\in\RR^m}\Bigl\{f(\bx)+h(\by)+\langle\blmbd,\calA(\bx)-\by\rangle
    +\frac{\rho}{2}\|\calA(\bx)-\by\|^2\Bigr\}.
    \label{eq:reduced-al}
\end{equation}
Its minimizer and the associated multiplier are
\begin{equation}
    \by(\bx,\blmbd):=\prox_{h/\rho}\Bigl(\calA(\bx)+\frac{\blmbd}{\rho}\Bigr),
    \qquad
    \bq(\bx,\blmbd):=\blmbd+\rho\bigl[\calA(\bx)-\by(\bx,\blmbd)\bigr],
    \label{eq:prox-multiplier}
\end{equation}
and \(\bq(\bx,\blmbd)\in\partial h(\by(\bx,\blmbd))\).
We study the smoothed perturbed linearized augmented Lagrangian method
(SPLALM), which performs alternating gradient descent--ascent steps on the
smoothed perturbed potential function
\begin{equation}
    \mathcal P_s(\bx,\blmbd;\bz,\bmu)
    :=\calD_\rho(\bx,\blmbd)+\frac{1}{2\gamma}\|\bx-\bz\|^2
    -\frac{1}{2s}\|\blmbd-\bmu\|^2 .
    \label{eq:potential}
\end{equation}
The proximal term smooths the primal problem, the term
\(-\frac{1}{2s}\|\blmbd-\bmu\|^2\) perturbs the dual problem, and
\(\bz\in\RR^n\) and \(\bmu\in\RR^m\) are the primal and dual centers.

\begin{definition}[SPLALM]
\label{def:splalm}
Let \(\theta=(\rho,s,\gamma,\eta,\tau)\) satisfy \(\rho>0\),
\(0<\eta\le\gamma<\infty\), and \(0<\tau\le s<\infty\), and let
\((\alpha,\beta)\in[0,1]^2\). The method \(\mathrm{SPLALM}_{\theta,\alpha,\beta}\)
starts from \((\bx^0,\blmbd^0,\bz^0,\bmu^0)=\zero\) and, for \(k\ge0\), sets
\begin{equation}
\begin{aligned}
    \bx^{k+1}&=\bx^k-\eta\,\nabla_{\bx}\mathcal P_s(\bx^k,\blmbd^k;\bz^k,\bmu^k),\\
    \blmbd^{k+1}&=\blmbd^k+\tau\,\nabla_{\blmbd}\mathcal P_s(\bx^{k+1},\blmbd^k;\bz^k,\bmu^k),\\
    \bz^{k+1}&=\bz^k+\alpha(\bx^{k+1}-\bz^k),\\
    \bmu^{k+1}&=\bmu^k+\beta(\blmbd^{k+1}-\bmu^k).
\end{aligned}
    \label{eq:splalm-updates}
\end{equation}
At iteration \(k\), it outputs
\((\bx^{k+1},\by^{k+1},\bq^k)
:=(\bx^{k+1},\by(\bx^{k+1},\blmbd^k),\bq(\bx^{k+1},\blmbd^k))\).
\end{definition}

With first-order information cached between iterations, \(K\) iterations and
their outputs use \(K+1\) calls to \(\mathcal O_{f,\calA}\) and \(2K\) calls to
\(\mathcal O_h\). Thus the total oracle count is \(3K+1\), and iteration
and oracle complexity agree up to a factor of three, apart from one initial
first-order oracle call.

\begin{definition}[SPLALM family]
\label{def:algorithm-family}
For \(\theta\) as in Definition~\ref{def:splalm},
\(\calC_\theta:=\{\mathrm{SPLALM}_{\theta,\alpha,\beta}:(\alpha,\beta)\in[0,1]^2\}\).
\end{definition}

\begin{remark}[Initialization]
\label{rem:lb-initialization}
Every method in \(\calC_\theta\) starts from zero on every instance.
Theorem~\ref{thm:lower-bound} constructs a hard instance for each
\((\theta,\alpha,\beta,\epsilon)\) while keeping this initialization fixed.
Thus, the lower bound holds for the prescribed starting point.
The constructions in~\cite[Section~5.1.1]{li2026smoothing}
and~\cite[Appendix~C.1]{mahdavinia2022tight} choose both the hard function
and the initialization.
\end{remark}

\begin{definition}[Worst-case iteration complexity]
\label{def:complexity-measure}
Let \(\epsilon>0\) and \(I\in\mathcal I(\mathbf C,\Delta)\), and let
\((\bx^{k+1},\by^{k+1},\bq^k)\) be the outputs of
\(\mathrm{SPLALM}_{\theta,\alpha,\beta}\) on \(I\). Define
\[
\begin{aligned}
    T^{\mathrm{KKT}}_\epsilon(\theta,\alpha,\beta;I)
    &:=\inf\Bigl\{K\ge1:\ \min_{0\le k<K}
        \mathcal R_{\mathrm{KKT}}(\bx^{k+1},\by^{k+1},\bq^k)\le\epsilon\Bigr\},\\
    T^{\mathrm{GS}}_\epsilon(\theta,\alpha,\beta;I)
    &:=\inf\Bigl\{K\ge1:\ \min_{0\le k<K}
        \calG(\bx^{k+1},\bq^k)\le\epsilon\Bigr\},
\end{aligned}
\]
with \(\inf\varnothing=+\infty\). For \(\bullet\in\{\mathrm{KKT},\mathrm{GS}\}\),
the worst-case iteration complexity of \(\calC_\theta\) over
\(\mathcal I(\mathbf C,\Delta)\) is
\[
    T^{\bullet}_\epsilon\bigl(\calC_\theta,\mathcal I(\mathbf C,\Delta)\bigr)
    :=\inf_{(\alpha,\beta)\in[0,1]^2}\
    \sup_{I\in\mathcal I(\mathbf C,\Delta)}\
    T^{\bullet}_\epsilon(\theta,\alpha,\beta;I).
\]
\end{definition}

The infimum precedes the supremum, so \((\alpha,\beta)\) may depend on
\(\epsilon\), \(\theta\), and \((\mathbf C,\Delta)\), but not on the instance.
At least one of the first \(K\) outputs meets the tolerance if and only if
\(K\ge T^{\bullet}_\epsilon\); hence the definition covers every rule that
returns one of these outputs, including the best-so-far rule. By
Proposition~\ref{prop:kkt-gs},
\(\calG(\bx^{k+1},\bq^k)\le\mathcal R_{\mathrm{KKT}}(\bx^{k+1},\by^{k+1},\bq^k)\)
for every \(k\), so
\(T^{\mathrm{GS}}_\epsilon(\calC_\theta,\mathcal I)
\le T^{\mathrm{KKT}}_\epsilon(\calC_\theta,\mathcal I)\). An upper bound for
the KKT measure therefore yields the same upper bound for the GS measure, and
a lower bound for the GS measure yields the same lower bound for the KKT
measure.

\begin{remark}[Family complexity and oracle complexity]
\label{rem:family-not-oracle}
Let \(\mathcal A_{\det}\) be the class of deterministic methods whose iterates
are arbitrary maps of all previous responses of \(\mathcal O_{f,\calA}\) and
\(\mathcal O_h\), defined as in~\cite[Section~2.2]{carmon2020lower}
and~\cite[Definition~2.5]{pan2026optimal} with the oracles of
Definition~\ref{def:oracle-model}. Then \(\calC_\theta\subset\mathcal A_{\det}\), and
\[
    T_\epsilon(\mathcal A_{\det},\mathcal I)
    =\inf_{A\in\mathcal A_{\det}}\sup_{I\in\mathcal I}T_\epsilon(A;I)
    \le\inf_{(\alpha,\beta)\in[0,1]^2}\sup_{I\in\mathcal I}
    T_\epsilon(\theta,\alpha,\beta;I)
    =T_\epsilon(\calC_\theta,\mathcal I).
\]
\begin{enumerate}[label=(\roman*)]
\item (Scope) A lower bound on \(T_\epsilon(\calC_\theta,\mathcal I)\) does not
  bound \(T_\epsilon(\mathcal A_{\det},\mathcal I)\); methods outside
  \(\calC_\theta\) may be faster. When \(h\) is the support function of a
  compact convex set, \eqref{eq:ncl} is a smooth nonconvex--concave problem,
  and the \(O(\epsilon^{-5/2})\) GS bound of~\cite{zhang2026matching} and the
  \(O(\epsilon^{-3})\) OS bound of~\cite{pan2026optimal} are attained by
  methods outside \(\calC_\theta\).
\item (Uniformity) The infimum ranges over all \((\alpha,\beta)\in[0,1]^2\),
  including the endpoint cases in Table~\ref{tab:special-cases}, and
  Theorem~\ref{thm:lower-bound} holds for all \((\rho,s)\) with
  \(\epsilon(\rho+s)\le c_1\). A lower bound on
  \(T_\epsilon(\calC_\theta,\mathcal I)\) therefore cannot be removed by
  tuning the center stepsizes; it is governed by the effective penalty
  \(\rho+s\) of Lemma~\ref{lem:effective-penalty}.
\item (Transfer) Conversely, every lower bound on
  \(T_\epsilon(\mathcal A_{\det},\mathcal I)\) for the same problem class and
  measure also holds for \(\calC_\theta\).
\end{enumerate}
\end{remark}

\subsection{Effective penalty and special cases}
\label{subsec:effective-penalty}

The reduced augmented Lagrangian is a dual proximal smoothing of the saddle
function \(\calL(\bx,\bq)=f(\bx)+\langle\calA(\bx),\bq\rangle-h^*(\bq)\)
in~\eqref{eq:ncl}~\cite{rockafellar1976augmented}. The center of this
smoothing is the multiplier, and its strength is the penalty. Adding the dual
regularizer of~\eqref{eq:potential} keeps this form; it only raises the
penalty and moves the center.

\begin{lemma}[Effective penalty]
\label{lem:effective-penalty}
Suppose that Assumption~\ref{ass:h} holds. For every \(\sigma>0\),
\begin{equation}
    \calD_\sigma(\bx,\blmbd)
    =\max_{\bq\in\RR^m}\Bigl\{\calL(\bx,\bq)-\frac{1}{2\sigma}\|\bq-\blmbd\|^2\Bigr\},
    \label{eq:dual-smoothing}
\end{equation}
and consequently
\begin{equation}
    \max_{\blmbd\in\RR^m}\mathcal P_s(\bx,\blmbd;\bz,\bmu)
    =\calD_{\rho+s}(\bx,\bmu)+\frac{1}{2\gamma}\|\bx-\bz\|^2 .
    \label{eq:effective-penalty}
\end{equation}
\end{lemma}

\begin{proof}
Write \(\bu=\calA(\bx)\) and \(h(\by)=\max_{\bq}\{\langle\bq,\by\rangle-h^*(\bq)\}\).
Since \(\dom h^*\) lies in a ball by Assumption~\ref{ass:h}, Sion's minimax
theorem allows the minimum over \(\by\) and the maximum over \(\bq\) in
\(\min_{\by}\max_{\bq}\{\langle\bq-\blmbd,\by\rangle-h^*(\bq)
+\frac{\sigma}{2}\|\bu-\by\|^2\}\) to be interchanged. The inner minimizer is
\(\by=\bu-(\bq-\blmbd)/\sigma\), which gives~\eqref{eq:dual-smoothing}.
Applying~\eqref{eq:dual-smoothing} with \(\sigma=\rho\) and exchanging the two
maximizations,
\[
    \max_{\blmbd}\mathcal P_s
    =\max_{\bq}\Bigl\{\calL(\bx,\bq)-\min_{\blmbd}\Bigl[\frac{\|\bq-\blmbd\|^2}{2\rho}
    +\frac{\|\blmbd-\bmu\|^2}{2s}\Bigr]\Bigr\}+\frac{\|\bx-\bz\|^2}{2\gamma},
\]
and the inner minimum equals \(\|\bq-\bmu\|^2/(2(\rho+s))\).
Then~\eqref{eq:dual-smoothing} with \(\sigma=\rho+s\)
gives~\eqref{eq:effective-penalty}.
\end{proof}

Lemma~\ref{lem:effective-penalty} shows that every member of
\(\calC_\theta\) has the effective penalty \(\rho+s\), whatever the center
stepsizes are. The center stepsize \(\beta\) only places the dual center
between two classical choices, and \(\alpha\) decides whether the primal
variable is smoothed.

\emph{Zero center} (\(\beta=0\)). Then \(\bmu^k=\zero\), and
\(\calD_{\rho+s}(\bx,\zero)=f(\bx)+\min_{\by}\{h(\by)
+\frac{\rho+s}{2}\|\calA(\bx)-\by\|^2\}\) is the quadratic penalty function
with parameter \(\rho+s\). The dual step carries the perturbation
\(-\blmbd^k/s\), as in perturbed GDA~\cite{lu2020hybrid,xu2023unified} and
perturbed smoothed GDA~\cite{li2026smoothing}.

\emph{Current multiplier as center} (\(\beta=1\)). Then
\(\bmu^k=\blmbd^k\), the gradient of the dual regularizer vanishes at every
dual step, and
\[
    \blmbd^{k+1}=\Bigl(1-\frac{\tau}{\rho}\Bigr)\blmbd^k+\frac{\tau}{\rho}\,\bq^k,
    \qquad
    \bq^k=\prox_{\rho h^*}\bigl(\blmbd^k+\rho\calA(\bx^{k+1})\bigr).
\]
For \(\tau=\rho\), this is the multiplier update of the augmented Lagrangian
method~\cite{hestenes1969multiplier,powell1969method,rockafellar1973multiplier},
that is, a proximal ascent step on \(\calL(\bx^{k+1},\cdot)\)
by~\eqref{eq:dual-smoothing}. The iterates do not depend on \(s\), so this
choice also reproduces the case without dual regularizer, \(s=\infty\).

\emph{Moving center} (\(0<\beta<1\)). The dual center is an exponentially
weighted average of past multipliers, and the method has the structure of doubly smoothed
GDA~\cite{zheng2023universal}.

Table~\ref{tab:special-cases} collects these cases. The identities
\(\bmu^k=\zero\), \(\bmu^k=\blmbd^k\), and \(\bz^k=\bx^k\) for
\(\beta=0\), \(\beta=1\), and \(\alpha=1\), respectively, follow by induction
from the zero initialization. Every entry refers to the cited update structure
applied to \(\calD_\rho\) with \(\rho>0\); methods acting on \(\calL\) itself
are compared in Section~\ref{subsec:gda-embedding}.

\begin{table}[t]
\centering
\small
\caption{Special cases of \(\calC_\theta\). All methods act on
\(\calD_\rho\) with \(\rho>0\) and have the effective penalty \(\rho+s\).}
\label{tab:special-cases}
\begin{tabular}{@{}llll@{}}
\toprule
\(\beta\) & Dual center & \(\alpha=1\) & \(0<\alpha<1\) \\
\midrule
\(0\) & \(\zero\)
  & perturbed GDA~\cite{lu2020hybrid,xu2023unified}
  & perturbed smoothed GDA~\cite{li2026smoothing} \\
\((0,1)\) & \(\sum_{j=1}^{k}\beta(1-\beta)^{k-j}\blmbd^j\)
  & ---
  & DS-GDA~\cite{zheng2023universal}, \(r_1=1/\gamma\), \(r_2=1/s\) \\
\(1\) & \(\blmbd^k\)
  & linearized ALM (\(\tau=\rho\))
  & smoothed GDA~\cite{zhang2020smoothed} \\
\bottomrule
\end{tabular}
\end{table}

